%=====================================================================
\chapter{Selecting Rows, and the Trouble with NULL}\label{ch:select}
%=====================================================================
\locator{4}{Part IV --- Asking Questions}

\begin{outcomes}
By the end of this chapter you will be able to:
\begin{tight}
  \item \textbf{Write} a \sqlkw{SELECT} with \sqlkw{WHERE}, \sqlkw{ORDER BY}
        and \sqlkw{LIMIT}, and \textbf{state} the order in which its clauses
        run.
  \item \textbf{Use} \sqlkw{LIKE}, \sqlkw{IN}, \sqlkw{BETWEEN} and
        \sqlkw{CASE} correctly, including at their boundaries.
  \item \textbf{Evaluate} a condition in three-valued logic, and
        \textbf{predict} which rows a condition involving null will return.
  \item \textbf{Explain} why \sqlkw{NOT IN} can return nothing, and
        \textbf{replace} it with a form that cannot.
  \item \textbf{Account for} every row of a table when a condition and its
        negation together do not cover it.
\end{tight}
\end{outcomes}

\begin{prereq}
\chref{relational}: what a null is. \chref{joins} for the 73 sections a null
foreign key dropped. A seeded \texttt{university.db}; every number in this
chapter is reproduced by its laboratory.
\end{prereq}

\begin{keyterms}
\sqlkw{SELECT} $\bullet$ \sqlkw{WHERE} $\bullet$ \sqlkw{ORDER BY} $\bullet$
\sqlkw{LIMIT} $\bullet$ \sqlkw{OFFSET} $\bullet$
\sqlkw{LIKE} $\bullet$ wildcard $\bullet$ \sqlkw{IN} $\bullet$
\sqlkw{BETWEEN} $\bullet$ \sqlkw{CASE} $\bullet$ three-valued logic $\bullet$
\unk{} $\bullet$ \sqlkw{IS NULL} $\bullet$ \sqlkw{COALESCE}
$\bullet$ \sqlkw{NULLIF} $\bullet$ \sqlkw{NOT IN} $\bullet$ \sqlkw{NOT EXISTS}
\end{keyterms}

\section{The Statement, and the Order It Runs In}

\begin{definitionbox}[\sqlkw{SELECT}, in execution order]
\begin{lstlisting}[language=sql]
SELECT   name, semester          -- 4. keep these columns
FROM     student                 -- 1. start from this relation
WHERE    programme = 'BSIT'      -- 2. keep these rows
ORDER BY semester DESC, name     -- 5. sort what is left
LIMIT    10 OFFSET 20;           -- 6. then take a slice
\end{lstlisting}

\smallskip
Step 3, between \sqlkw{WHERE} and \sqlkw{SELECT}, is \sqlkw{GROUP BY}, and
\chref{grouping} adds it. The order matters in one practical way: a name given
in the \sqlkw{SELECT} list with \sqlkw{AS} \emph{does not exist yet} when
\sqlkw{WHERE} runs, so it cannot be used there --- but it can be used in
\sqlkw{ORDER BY}, which runs after.
\end{definitionbox}

\begin{examplebox}[The four predicates, and where each one surprises]
\begin{tabular}{@{}L{4.6cm}L{9.4cm}@{}}
\texttt{name LIKE 'A\%'} & the \term{wildcards}: \texttt{\%} is any run of characters,
  \texttt{\_} exactly one. In SQLite \sqlkw{LIKE} is \textbf{case-insensitive}
  for ASCII letters, so this also matches \texttt{'ayesha'}. \\
\texttt{programme IN ('BSIT','BSCS')} & shorthand for two \sqlkw{OR}s. Fine
  --- until the list contains a null, below. \\
\texttt{semester BETWEEN 2 AND 4} & \textbf{inclusive at both ends}: 2, 3 and
  4. \\
\texttt{CASE WHEN g='A' THEN 4 \dots\ END} & evaluated top to bottom; the
  first true branch wins; with no \sqlkw{ELSE}, the result is null. \\
\end{tabular}

\smallskip
\textbf{And one that is not a predicate at all.} \texttt{LIMIT 10} without
\sqlkw{ORDER BY} returns \emph{some} ten rows. Which ten depends on the plan,
and \chref{relational} said the plan is allowed to change. A ``top ten'' with
no \sqlkw{ORDER BY} is ten rows chosen by the optimiser.
\end{examplebox}

\section{Three-Valued Logic}

Every condition so far has been true or false. As soon as a null is involved
there is a third answer, and it is the source of most quietly wrong results in
SQL.

\begin{definitionbox}[\unk{}]
Any comparison with a null --- \texttt{=}, \texttt{<>}, \texttt{<},
\texttt{>}, even \texttt{NULL = NULL} --- evaluates to \unk{}, not to
true or false. It means ``the value is missing, so I cannot say''.

\smallskip
\sqlkw{WHERE} keeps a row \textbf{only if its condition is true}. False and
\unk{} are both discarded.

\smallskip
\sqlkw{IS NULL} and \sqlkw{IS NOT NULL} are the only tests that return true
or false for a null. They exist because \texttt{= NULL} cannot work.
\end{definitionbox}

\begin{conceptbox}[The truth tables, and the one row to remember]
\begin{tabular}{@{}L{1.9cm}L{1.9cm}|L{1.9cm}L{1.9cm}|L{2.2cm}L{2.4cm}@{}}
\textbf{\sqlkw{AND}} & & \textbf{\sqlkw{OR}} & & \textbf{\sqlkw{NOT}} & \\
T \sqlkw{AND} U & = U & T \sqlkw{OR} U & = \textbf{T} & \sqlkw{NOT} T & = F \\
F \sqlkw{AND} U & = \textbf{F} & F \sqlkw{OR} U & = U & \sqlkw{NOT} F & = T \\
U \sqlkw{AND} U & = U & U \sqlkw{OR} U & = U & \sqlkw{NOT} U & = \textbf{U} \\
\end{tabular}

\smallskip
Read the bold entries. \sqlkw{AND} with a false is false whatever the other
side; \sqlkw{OR} with a true is true whatever the other side --- those are the
cases where a missing value genuinely does not matter.

\smallskip
\textbf{The row to remember is the last one: \sqlkw{NOT} \unk{} is
\unk{}.} Negating a condition does \emph{not} give you the other
rows. It gives you the other rows \emph{that are not null}. A condition and its
negation together do not cover a table that contains nulls --- and the next box
counts exactly how many rows fall through the gap.
\end{conceptbox}

\begin{worked}[Grade A, not grade A, and the 19{,}992 rows that are neither]
The seeded \textit{enrolment} table has 60{,}037 rows. Three queries:

\medskip
\begin{tabular}{@{}L{7.2cm}R{2.6cm}@{}}
\texttt{WHERE grade = 'A'}       & 7{,}268 \\
\texttt{WHERE grade <> 'A'}      & 32{,}777 \\
\texttt{WHERE NOT (grade = 'A')} & 32{,}777 \\
\end{tabular}

\medskip
\textbf{Step 1 --- the sum.} $7{,}268 + 32{,}777 = 40{,}045$. The table has
60{,}037. \textbf{19{,}992 rows are in neither answer.}

\medskip
\textbf{Step 2 --- which rows.} The FALL2026 enrolments, whose semester is not
over: \texttt{grade} is null. For those rows \texttt{grade = 'A'} is
\unk{}, so is \texttt{grade <> 'A'}, and so --- by the last row of
the truth table --- is \sqlkw{NOT} \texttt{(grade = 'A')}. Every one of the
three conditions discards them.

\medskip
\textbf{Step 3 --- the check.}
\begin{lstlisting}[language=sql]
SELECT COUNT(*) FROM enrolment WHERE grade IS NULL;    -- 19992
\end{lstlisting}
$7{,}268 + 32{,}777 + 19{,}992 = 60{,}037$. Every row accounted for.

\medskip
\textbf{Step 4 --- what goes wrong in practice.} A report of ``students who
did not get an A'' built on \texttt{grade <> 'A'} omits a third of all
enrolments without a word. Whether it \emph{should} include them is a question
for the examinations office --- an ungraded course is not a failure --- but
the query decided it silently, and decided it the way most people would not.

\medskip
\textbf{The rule.} When a column can be null, write the null case
explicitly, whichever way you want it:
\begin{lstlisting}[language=sql]
WHERE grade <> 'A' OR grade IS NULL     -- include ungraded
WHERE grade <> 'A'                      -- exclude them, and say so
\end{lstlisting}
\end{worked}

\dsfig{%
\begin{tikzpicture}[
  seg/.style={draw=white, line width=1pt, minimum height=0.9cm,
              font=\scriptsize\bfseries, text=white, inner sep=0pt},
  br/.style={draw=#1, line width=0.9pt,
             decorate, decoration={brace, amplitude=5pt}},
  lb/.style={font=\scriptsize, align=center, text=#1}
]
  % 60,037 rows drawn to scale across 13 cm. The widths are computed
  % here, not with pgfmath: the first draft used macros named a, b and c
  % -- which shadow TeX's own accent commands -- and a 13/60037 scale
  % factor, and three 'Dimension too large' errors followed.
  %   13 x 7268/60037 = 1.5738    13 x 32777/60037 = 7.0973
  %   13 x 19992/60037 = 4.3289
  \def\wA{1.5738}\def\wN{7.0973}\def\wU{4.3289}
  \node[seg, fill=greenhl, minimum width=\wA cm, anchor=west]
       (A) at (0,0) {A};
  \node[seg, fill=secondary, minimum width=\wN cm, anchor=west]
       (N) at (\wA,0) {B, C, D, F};
  \node[seg, fill=neutral!70, minimum width=\wU cm, anchor=west]
       (U) at (\wA+\wN,0) {NULL --- not yet graded};

  \draw[br=greenhl] (0,0.55) -- (\wA,0.55);
  \node[lb=greenhl!60!black] at (\wA/2,1.05)
       {\texttt{grade = 'A'}\\7{,}268};
  \draw[br=secondary] (\wA,0.55) -- (\wA+\wN,0.55);
  \node[lb=secondary!70!black] at (\wA+\wN/2,1.05)
       {\texttt{grade <> 'A'} \ and \ \texttt{NOT (grade = 'A')}\\32{,}777};
  \draw[br=accent] (\wA+\wN+\wU,-0.55) -- (\wA+\wN,-0.55);
  \node[lb=accent] at (\wA+\wN+\wU/2,-1.15)
       {\textbf{caught by neither}\\19{,}992 --- \sqlkw{IS NULL} only};
  \node[lb=neutral, anchor=east] at (-0.15,0) {60{,}037\\rows};
\end{tikzpicture}}%
{The seeded \textit{enrolment} table, drawn to scale. A condition and its
negation cover the green and the blue. The grey third is \unk{} to
both, and \sqlkw{WHERE} discards it either way.}{nullgap}

\section{\sqlkw{NOT IN}, Which Can Return Nothing}

\begin{alertbox}[Three new instructors, and a query that cannot find them]
Hire three instructors who are not yet teaching anything. Then ask which
instructors teach no section --- the obvious way:

\begin{lstlisting}[language=sql]
SELECT COUNT(*) FROM instructor
 WHERE emp_id NOT IN (SELECT emp_id FROM section);
\end{lstlisting}

On the seeded database, with the three added, this returns \textbf{0}.

\begin{lstlisting}[language=sql]
SELECT COUNT(*) FROM instructor i
 WHERE NOT EXISTS (SELECT 1 FROM section s WHERE s.emp_id = i.emp_id);
\end{lstlisting}

This returns \textbf{3}.

\medskip
\textbf{Why.} \texttt{x NOT IN (a, b, c)} means \texttt{x <> a AND x <> b AND
x <> c}. The sub-query's list contains the 73 nulls from the unassigned
sections, so the expansion contains \texttt{x <> NULL} --- which is
\unk{}. By the truth table, \texttt{T AND U = U}. For every
instructor, including the three new ones, the whole condition is at best
\unk{}, and \sqlkw{WHERE} discards \unk{}. \textbf{One null
in the list empties the result.}

\medskip
\textbf{The fix is either of these:}
\begin{tight}
  \item use \sqlkw{NOT EXISTS}, which tests for the \emph{existence} of a row
        rather than comparing with a value, and is unaffected by nulls; or
  \item filter the nulls out of the list:
        \texttt{\dots\ NOT IN (SELECT emp\_id FROM section WHERE emp\_id IS NOT
        NULL)}, which also returns 3.
\end{tight}

\smallskip
\textbf{This course uses \sqlkw{NOT EXISTS}.} The second fix works, but it has
to be remembered every time, and it will be forgotten the first time the
column becomes nullable --- which is a schema change somebody else will make.
\chref{subqueries} measures the two forms against each other.
\end{alertbox}

\begin{examplebox}[\sqlkw{COALESCE} and \sqlkw{NULLIF}: deciding what null means]
\sqlkw{COALESCE}(a, b, c) returns the first argument that is not null. It is
how you \emph{decide}, in one place, what a null should be treated as:

\begin{lstlisting}[language=sql]
SELECT roll_no, COALESCE(cnic, 'not on file') AS cnic FROM student;
SELECT COALESCE(grade, 'in progress') AS grade FROM enrolment;
\end{lstlisting}

\sqlkw{NULLIF}(a, b) is the reverse: it returns null if $a = b$, otherwise
$a$. Its main use is cleaning up \chref{relational}'s magic values ---
\texttt{NULLIF(cnic, 'N/A')} turns an invented marker back into the null it
should have been.

\smallskip
\textbf{Use \sqlkw{COALESCE} for display, not for arithmetic.} Writing
\texttt{AVG(COALESCE(score, 0))} decides that a missing mark is a zero, which
is \chref{ddl}'s 20\% error made on purpose.
\end{examplebox}

\begin{conceptbox}[Where nulls go in \sqlkw{ORDER BY}]
SQLite sorts nulls \textbf{first} in ascending order and last in descending,
treating null as smaller than every value. Other systems differ: PostgreSQL
puts them last in ascending order. SQLite since 3.30 accepts the standard
clause that makes it explicit:

\begin{lstlisting}[language=sql]
SELECT roll_no, grade FROM enrolment ORDER BY grade NULLS LAST;
\end{lstlisting}

Write \sqlkw{NULLS FIRST} or \sqlkw{NULLS LAST} whenever the column can be
null and the order matters, because ``what this database happens to do'' is
not the same on the next one.
\end{conceptbox}

\begin{pitfall}
\begin{tight}
  \item \textbf{\texttt{WHERE x = NULL}.} Always \unk{}, so always
        empty. Use \sqlkw{IS NULL}.
  \item \textbf{Assuming a condition and its negation cover the table.} They
        miss the nulls --- 19{,}992 of 60{,}037 enrolments.
  \item \textbf{\sqlkw{NOT IN} against a nullable column.} One null in the
        list and the result is empty. Use \sqlkw{NOT EXISTS}.
  \item \textbf{Using a column alias in \sqlkw{WHERE}.} It does not exist yet.
  \item \textbf{\sqlkw{LIMIT} without \sqlkw{ORDER BY}.} Some rows, chosen by
        the plan.
  \item \textbf{Forgetting \sqlkw{BETWEEN} is inclusive.} Both ends.
  \item \textbf{Expecting \sqlkw{LIKE} to be case-sensitive.} In SQLite it is
        not, for ASCII.
  \item \textbf{\sqlkw{COALESCE} to zero inside an average.} A decision that a
        missing mark is a fail, made silently.
\end{tight}
\end{pitfall}

\begin{lab}[Lose 19{,}992 rows and three instructors, and find them all]
\begin{lstlisting}[language=bash]
# --- 1. A fresh seeded database --------------------------------------
python code/schema/seed.py university.db
sqlite3 university.db
\end{lstlisting}
\begin{lstlisting}[language=sql]
-- --- 2. A condition and its negation ----------------------------
.headers on
.mode box
SELECT COUNT(*) FROM enrolment WHERE grade = 'A';
SELECT COUNT(*) FROM enrolment WHERE grade <> 'A';
SELECT COUNT(*) FROM enrolment;
-- Add the first two. Compare with the third. Explain the gap.

-- --- 3. Account for every row -----------------------------------
SELECT COUNT(*) FROM enrolment WHERE grade IS NULL;
SELECT COUNT(*) FROM enrolment WHERE grade = NULL;
-- Two queries that look alike. One finds the gap; one finds nothing.

-- --- 4. Hire three people who teach nothing yet -----------------
PRAGMA foreign_keys = ON;
INSERT INTO instructor VALUES
  ('E-9001','New Hire One','Lecturer','CS'),
  ('E-9002','New Hire Two','Lecturer','IT'),
  ('E-9003','New Hire Three','Lecturer','SE');

-- --- 5. Look for them with NOT IN -------------------------------
SELECT COUNT(*) FROM instructor
 WHERE emp_id NOT IN (SELECT emp_id FROM section);
-- There are three. How many did it find?

-- --- 6. Look for them with NOT EXISTS ---------------------------
SELECT name FROM instructor i
 WHERE NOT EXISTS (SELECT 1 FROM section s WHERE s.emp_id = i.emp_id);

-- --- 7. Prove the nulls are to blame ----------------------------
SELECT COUNT(*) FROM section WHERE emp_id IS NULL;
SELECT COUNT(*) FROM instructor
 WHERE emp_id NOT IN (SELECT emp_id FROM section
                       WHERE emp_id IS NOT NULL);

-- --- 8. Decide what null means, once ----------------------------
SELECT COALESCE(grade, 'in progress') AS g, COUNT(*)
  FROM enrolment GROUP BY g ORDER BY g;
SELECT roll_no, grade FROM enrolment ORDER BY grade NULLS LAST LIMIT 3;
\end{lstlisting}

\smallskip
\textbf{What to notice.} Step 5 against step 6 is the chapter in one
comparison: the same question, two spellings, 0 against 3. Neither raised an
error, and the wrong one is the one most people write first. If you remember
one rule from Part~IV, make it this: \textbf{never \sqlkw{NOT IN} a column
that can be null.}
\end{lab}

\begin{checkpoint}
\begin{tightnum}
  \item \texttt{grade = 'A'} returns 7{,}268 rows and \texttt{grade <> 'A'}
        returns 32{,}777, on a table of 60{,}037. Account for the remaining
        rows using the truth table, and write the single query that returns
        every row that is \emph{not} an A, including the ungraded ones.
  \item Explain, step by step from the expansion of \sqlkw{NOT IN} into
        \sqlkw{AND}s, why one null in the sub-query's result makes the outer
        query return no rows at all.
  \item A colleague writes \texttt{AVG(COALESCE(score, 0))} ``so that the
        average includes everybody''. Say what decision this makes about
        missing marks, show with a three-row example how it changes the
        average, and say what you would write instead and why.
\end{tightnum}
\end{checkpoint}

\begin{chaptersummary}
\begin{tight}
  \item \sqlkw{SELECT} runs \sqlkw{FROM}, \sqlkw{WHERE}, \sqlkw{GROUP BY},
        \sqlkw{SELECT}, \sqlkw{ORDER BY}, \sqlkw{LIMIT}. An alias does not
        exist in \sqlkw{WHERE}; it does in \sqlkw{ORDER BY}.
  \item \sqlkw{BETWEEN} is inclusive; SQLite's \sqlkw{LIKE} is
        case-insensitive for ASCII; \sqlkw{LIMIT} without \sqlkw{ORDER BY}
        returns rows the plan chose.
  \item Any comparison with a null is \textbf{\unk{}}.
        \sqlkw{WHERE} keeps only \emph{true}, so \unk{} rows vanish.
  \item \textbf{\sqlkw{NOT} \unk{} is \unk{}}: a condition
        and its negation do not cover a table with nulls.
  \item \textbf{On the seeded data:} \texttt{= 'A'} found 7{,}268,
        \texttt{<> 'A'} found 32{,}777, and \textbf{19{,}992 ungraded rows were
        in neither}.
  \item \texttt{= NULL} is always empty. \sqlkw{IS NULL} is the only test.
  \item \textbf{\sqlkw{NOT IN} returned 0 where \sqlkw{NOT EXISTS} returned 3},
        because one null in the list makes every comparison \unk{}.
        Never \sqlkw{NOT IN} a nullable column.
  \item \sqlkw{COALESCE} decides what null means; use it for display, never to
        turn a missing mark into a zero.
  \item SQLite sorts nulls first ascending; say \sqlkw{NULLS LAST} when it
        matters, because other systems differ.
\end{tight}
\end{chaptersummary}

\begin{reviewq}
  \item \textbf{(MCQ)} \texttt{NULL = NULL} evaluates to: (a)~true;
        (b)~false; (c)~\unk{}; (d)~an error.
  \item \textbf{(MCQ)} \texttt{semester BETWEEN 2 AND 4} includes:
        (a)~3 only; (b)~2 and 3; (c)~3 and 4; (d)~2, 3 and 4.
  \item \textbf{(Short)} Why can a column alias be used in \sqlkw{ORDER BY}
        but not in \sqlkw{WHERE}?
  \item \textbf{(Short)} Evaluate \texttt{F OR U}, \texttt{T AND U} and
        \sqlkw{NOT} \texttt{U}, and say which one causes the most errors in
        practice.
  \item \textbf{(Short)} Give the two fixes for a \sqlkw{NOT IN} against a
        nullable column, and say why this course prefers one of them.
  \item \textbf{(Applied)} Write a query listing every student with a label
        computed by \sqlkw{CASE}: \texttt{'senior'} for semesters 7--8,
        \texttt{'mid'} for 3--6, \texttt{'new'} for 1--2. Then say what your
        query does with a null semester, and change it so that a null is
        labelled \texttt{'unknown'} explicitly.
  \item \textbf{(Applied)} You are asked for ``students who have never failed
        a course''. Write it three ways --- with \sqlkw{NOT IN}, with
        \sqlkw{NOT EXISTS}, and with \sqlkw{EXCEPT} --- run all three on the
        seeded database, and explain any difference between their counts in
        terms of the null grades of FALL2026.
\end{reviewq}
